\documentclass[11pt,a4paper]{article}
\usepackage[spanish,es-noshorthands,es-nodecimaldot,es-nosectiondot,es-noindentfirst]{babel}
\usepackage{fontspec,amsmath,amssymb,mathtools,amsthm,mathrsfs}
\usepackage[a4paper,margin=26mm]{geometry}
\usepackage{xcolor,microtype,booktabs}
\usepackage[unicode,colorlinks=true,linkcolor={blue!45!black},urlcolor={blue!55!black}]{hyperref}
\hypersetup{pdftitle={Las tres estimaciones de llegada del lema B.1},
  pdfauthor={Dilan Steven Mejia},pdflang={es}}
\IfFontExistsTF{texgyrepagella-regular.otf}
 {\setmainfont{texgyrepagella}[Extension=.otf,UprightFont=*-regular,BoldFont=*-bold,
  ItalicFont=*-italic,BoldItalicFont=*-bolditalic]}{\setmainfont{Latin Modern Roman}}
\setlength{\emergencystretch}{3em}
\setcounter{tocdepth}{1}
\allowdisplaybreaks[2]
\newcommand{\E}{\mathbb E}
\newcommand{\EE}{\mathbb E}
\newcommand{\PP}{\mathbb P}
\newcommand{\TV}{\operatorname{TV}}
\newcommand{\per}{\operatorname{per}}
\newcommand{\op}{\mathrm{op}}
\newcommand{\diag}{\operatorname{diag}}
\newcommand{\tr}{\operatorname{tr}}
\newcommand{\pdet}{\operatorname{pdet}}
\newcommand{\Log}{\operatorname{Log}}
\newtheorem{proposicion}{Proposición}
\newtheorem{teorema}{Teorema}
\newtheorem{lema}{Lema}
\newcommand{\pruebaen}[1]{\hyperref[pieza:#1]{\textsf{#1}}}
\newcommand{\pieza}[1]{\begingroup
  \renewcommand{\section}[1]{\subsection{##1}\label{pieza:#1}%
    \noindent{\small\textsf{Fuente #1}}\par\nobreak\smallskip}
  \input{fuentes/#1}\endgroup}
\title{Las tres estimaciones de llegada\\del lema B.1}
\author{Dilan Steven Mejia}
\date{}
\begin{document}
\maketitle
\noindent Suplemento a \emph{La corrección de orden \(1/n\) al predictor canónico de las \(n\) reinas}.
Fijado un tablero de llegada, hay dos logaritmos que comparar. El primero mide
la probabilidad de evitar los ataques diagonales al colocar torres. El segundo es el cambio
del predictor al activar esos ataques. Las tres proposiciones de este cuaderno permiten
compararlos después de sustraer el aporte de cada cuadrado.
\tableofcontents
\smallskip
\noindent Las secciones 1--4 fijan las hipótesis y enuncian las estimaciones con su prueba.
Las secciones 5--7 desarrollan los cálculos empleados. Las claves enlazan con esos
desarrollos. El archivo \texttt{LEEME.md} identifica los antecedentes adicionales y
\texttt{ORIGEN.tsv} da la procedencia de cada pieza. En los originales, B y C designan
los artículos españoles [12] y [13] de la bibliografía del preprint, respectivamente.
\section{La clase en el hijo de tamaño \texorpdfstring{\(j\)}{j}}
Fijemos la banda de [13, §5.4], la entrada regular en \(M=\lceil\delta n\rceil\), la difusión
\(C\), la brecha \(\gamma\) y las constantes del coste exterior y de salida \(c_*>2\beta\).
En \(H_j\) la historia ha permanecido regular hasta el hijo de tamaño \(j=\lfloor bn\rfloor\)
inclusive. Cada línea diagonal tiene menos de \(\lceil\alpha j\rceil\) casillas disponibles.
Denotemos por \(A\) su máscara binaria. Se usan tres matrices en este mismo soporte.
El punto de reinas \(p\) cumple \(j\max p_e\le C\) y tiene brecha de incidencia de filas y columnas
al menos \(\gamma\). La masa \(p\) de cada diagonal es a lo sumo \(C\alpha\) por integralidad.

El punto estacionario de torres \(P\) satisface en el soporte
\(|P_e/p_e-1|\le K_T(C\alpha+1/j)\) por la contracción en el tangente de
\pruebaen{D04}. Con \(C\alpha<\epsilon_T\), \(K_TC\alpha\le1/4\) y \(K_T/j\le1/4\)
se obtiene \(p/2\le P\le 3p/2\). En particular \(j\max P_e\le 2C\) y
\(\epsilon_A:=\max_f\sum_{e\in f}P_e\le2C\alpha\) en ambas familias diagonales.
La forma de incidencia de \(P\) domina la mitad de la de \(p\) en el complemento
del modo nulo \((\mathbf1,-\mathbf1)\). Sus autovalores son \(1\pm\sigma_i(P)\).
Así \(\|P-J_j\|_{\rm op}\le1-\gamma_T\), con
\(J_j=\mathbf1\mathbf1^t/j\) y \(\gamma_T=\min(1,\gamma)/2\).

Sea \(D\) el escalado bistocástico exacto de \(A\). La estacionariedad da
\(P_e(1-P_e)=a_rb_c\) en el soporte. Balancear \(P\) con el campo
\(\log(1-P_e)=O(1/j)\) da, por el inverso del jacobiano de márgenes en el calibre de
\pruebaen{N02}, \(D_e/P_e=1+O(1/j)\). La unicidad del escalado lo identifica con \(D\).
El error tiene sumas absolutas por fila y por columna de orden \(1/j\). La desigualdad
\(\|D-P\|_2\le\sqrt{\|D-P\|_1\|D-P\|_\infty}\) da la misma cota en norma de operador.
Para \(j\) grande pueden reservarse \(j\max D_e\le4C\) y
\(\|D-J_j\|_2\le1-\gamma_T/2\). Hall, aplicado a \(D\), da \(\per A>0\).

\paragraph{Orden de elección.}
Los márgenes de clase fijan primero los radios de balanceo y la fracción máxima de
supresiones \(a_*\), y luego \(\varepsilon\) para el orden común
\(Q_*=\lfloor\varepsilon j\rfloor\). Para la norma exacta se reserva el factor exponencial
por sitio de las prescripciones largas. Se escogen después la fracción de corte de Taylor,
la de prescripciones cortas y el umbral \(\eta_F\) en \pruebaen{V03}.
Sea \(\eta=\min(\eta_F,\eta_\Lambda)\), con el umbral canónico de \pruebaen{N03}.
Independientemente, los mismos márgenes fijan los radios locales y la base combinatoria
\(C_1\) con alcance \(8\log n\). Fíjense \(0<a<1\) con
\(a\log(\max(2,C_1))<1/4\) y después \(\rho\ge2\) con \(a\log\rho>3\).
Se toma \(\alpha<\alpha_{\rm hist}\) para satisfacer las condiciones de torres anteriores y
\(2C\alpha<\eta/(2\rho^2)\), reduciéndola para los márgenes de carga del exterior en \(D\),
si se expresan en esa tabla. Para \(j\) grande, \(\rho^2(\epsilon_A+1/j)\le\eta\).

Se reduce \(\delta_*\) para ambas series geométricas de B.1, se elige el corte
\(0<b<\min(\delta/2,\delta_*)\) con su condición de prescripción y se hace \(n\) grande
a \(b\) fijo. La constante de prescripción procede de la prueba de [12, teorema 8.1]
en la clase histórica. En particular se fija antes de \(b\). Se obtiene
\(j\ge\lceil\sqrt n\rceil\), \(C'(4v+12)\le j^{1/3}\) y \(v^2/j\le\delta_*\)
para \(v\le8\log n\). Estas son condiciones sobre \(n_0(b)\).

\paragraph{Activación por sitios.}
En una evaluación \(Y\subseteq A\), se mantienen las restricciones de filas y columnas
y se activan los conflictos diagonales entre sitios de \(Y\). Sea \(Z(Y)\) la probabilidad
de evitarlos bajo el matching uniforme de \(A\), \(F(Y)=\log Z(Y)\) una vez demostrada
la positividad, y \(\Lambda(Y)\) el cambio del predictor respecto de torres.
Ponemos \(e=F-\Lambda\) y
\(\Delta_U f=\sum_{Y\subseteq U}(-1)^{|U|-|Y|}f(Y)\).
Los tres resultados se prueban sobre máscaras fijas. La probabilidad de torres evalúa
\(F\) dentro de cada máscara. La posterior integración de B.1 usa el indicador \(H_j\)
bajo la ley original de historias \(\nu\).

\section{Comparación de pesos en todos los menores de una cara}
\begin{proposicion}\label{prop:m}
Supónganse los márgenes de \(D\) de la sección 1. Fíjese \(U\) con
\(5\le v=|U|\le8\log n\). En todas las caras se toma el mismo exterior \(A_0=A\setminus U\).
Para cada \(I\subseteq U\) compatible con filas, columnas y ataques activos, el peso binario
\(W_I=\per(A_{0,-I})\) cuenta las configuraciones cuya intersección con \(U\) es exactamente
\(I\). Aquí \(-I\) suprime las filas y columnas ocupadas. Hay aproximantes positivos que
conservan los factores de escalado y el normalizador individual del menor, con
\[
 \left|\log(\widehat W_I/W_I)\right|\le \eta_v=(Cq/j)^q+Ce^{-c j},
 \qquad q=4v+12 .
\]
Los átomos prohibidos se omiten en ambas sumas. Si \(\widehat F_{U,q}\) usa esos pesos,
la suma de \(|\Delta_U(\widehat F_{U,q}-F)|\) en todo el intervalo es \(O(j^{-8})\).
\end{proposicion}
\begin{proof}
Al quitar los \(v\) sitios, la pérdida por margen es \(O(v/j)\). Al suprimir \(h=|I|\)
índices por lado, la variación de las medidas normalizadas es a lo sumo \(2h/j\) por lado.
El balanceo de \pruebaen{U02} invierte \(I-D^tD\) en columnas centradas y extiende la
solución a los índices de masa cero. En cada menor real de tamaño \(m=j-h\), la difusión
es uniforme y el cambio de la brecha cuesta \(O(\sqrt{h/j}+\sqrt v/j+(h+v)/j)\).
Así el escalado propio \(D'\) existe, es positivo en el soporte y
\(B=D'-J_m\) satisface \(\|B\|_2\le r'<1\). Hall aplicado al escalado da la
positividad de cada menor compatible, incluido el exterior vacío de marcas ocupadas.

Para \(R_B(z)=\sum_S z^{|S|}\prod_{e\in S}B_e\) sobre los matchings, desarrollar
\(\per(J_m+B)\) da \(m^m\per D'/m!=\mathbb E R_B(m/T)\), con \(T\sim\Gamma(m+1,1)\).
La integral de Wick de \pruebaen{V01} y la inversión sobre índices coincidentes en
\pruebaen{U03} dan \(S_k(y,u)=d(y)(1+\sum_{\ell=1}^k V_\ell(y,u))\), con
\(d(y)=\det(I-y^2 B^tB)^{-1/2}\) y \(|V_\ell|\le(C\ell|u|)^\ell\). En el disco auxiliar
\(|w|\le r_1>1\) se comparan primero los coeficientes hasta \(Q_*\) en \(R_B(yw)/d(yw)\).
El resto es \(O(e^{C\sqrt m}r_1^{-Q_*})=O(e^{-cj})\) para \(y\) cerca de 1.
Fuera de esa ventana, la cola gamma desplazada por \(\sum_i\sigma_i|Z_i|^2\), con los
valores singulares \(\sigma_i\) de \(B\), tiene coste ponderado \(O(e^{-cm+C\sqrt m})\)
por \pruebaen{V01}. En particular se domina el integrando completo en la cola.

Los momentos de \(m/T-1\) satisfacen la recurrencia de \pruebaen{U04}. Integrando el
Taylor de \(S_q(y,1/m)\) hasta \(2q+1\) se obtiene \(H_q(1/m)\). El momento par
\(2q+2\) acota el resto por \((C(q+1)/m)^{q+1}\). La suma de los excesos entre
\(q+1\) y \(Q_*\) tiene esta misma cota por razón geométrica. El orden \(Q_*\) se ha fijado
en el padre, mientras que \(1/m\) siempre corresponde al menor. Como \(H_q=d(1)+O(1/m)\)
y \(d(1)\) está separado de cero, la comparación vale en logaritmos. Escalado y factorial
se reponen en \(W_I\). En la normalización del predictor, el cambio de covolumen del
cofactor es \(\tfrac12\log(m/j)\) por \pruebaen{U07}. Equivale a usar el tamaño \(m\)
en el menor. La resta de dos errores logarítmicos y \(m\ge j/2\) se absorben en \(C\).

En cada cara la positividad permite sumar las dos comparaciones multiplicativas. El
logaritmo cuesta \(\eta_v\) y la alternancia cuesta \(2^v\eta_v\). El denominador basal
es el mismo. Por \pruebaen{A03}, hay a lo sumo \(j^{2v}\) soportes del grado \(v\).
La parte de potencia queda acotada por \(2^v j^{-8-2v/3}\), y la parte exponencial
total por \(\exp(O(\log^2 n)-cj)\). La suma es \(O(j^{-8})\).
\end{proof}

\section{Los tres exponentes en un soporte}
\begin{proposicion}\label{prop:s}
Bajo los márgenes de la sección 1, para \(5\le v\le8\log n\), sea \(c_s\) el número
de componentes del grafo bipartito de los sitios de \(U\) y sus fibras de las cuatro
familias. Sea \(\beta_U\) su número total de ciclos independientes, y
\(d_i=\max_{f\ni i}|f\cap U|\). Para el aproximante de la proposición~\ref{prop:m},
\[
 \left|\Delta_U(\widehat F_{U,q}-\Lambda)\right|
 \le C^v(Cv)^{c_s+1}j^{-p}\prod_{i\in U}d_i,\qquad
 p=\begin{cases}
 v+c_s+1 & \beta_U=0,\\
 v+c_s & \beta_U=1,\\
 v+c_s-1 & \beta_U\ge2 .
 \end{cases}
\]
La cota se aplica a los soportes pertinentes, en los que cada sitio participa en un
conflicto activable. Filas y columnas permanecen en el grafo aunque se apague una diagonal.
\end{proposicion}
\begin{proof}
Sea \(s=1/j\) en el punto físico y libérense los campos \(\zeta_i\) de cada sitio
después de extraer el mismo campo afín en todas las caras. El balanceo de \pruebaen{B01} y
\pruebaen{B02} admite las deformaciones de medida de tamaño \(O(v|s|)\).
Por \(\sum_{i\in f}1/d_i\le1\), la silla local y su determinante en \pruebaen{G02}
son holomorfos en \(|s|<\kappa/v,\ |\zeta_i|<\kappa'/d_i\). Las diferencias mixtas
de los pesos en \pruebaen{G03} y la inducción ocupado/vacío en \pruebaen{G04} dan
un logaritmo del conteo aproximado \(O(v)\) en esa caja. Para la inducción basta
\(\prod_D(1+|e^{a_D}-1|)\le\exp(\sum_D|a_D|)\), por la serie exponencial.
La rama del predictor coincide con la solución pequeña de \pruebaen{N02}.
El cofactor común de \pruebaen{U07} hace coincidir las caras en \(\zeta_i=0\).
El alternado \(H(s,\zeta)\) está pues acotado por \(C^v\) y es divisible por todos los
\(\zeta_i\).

Para contar potencias de \(s\), se toma un grupo por componente. Un monomio conexo de coste
\(w\), con interacciones de costes \(\ell_a\) que tocan \(h_a\) grupos distintos, ocupa
\(w+1-\sum_a(\ell_a+1-h_a)-L_I\) grupos, donde \(L_I\ge0\) es el número de ciclos del
grafo de incidencias del monomio. Los déficits son no negativos por las dos fórmulas del
exterior de \pruebaen{B01} y su desarrollo en \pruebaen{P03}. Los coeficientes sobre grupos
separados se factorizan antes de tomar logaritmos. Al alternar, un grupo ausente da cero.
Por ello desaparecen todos los costes inferiores a \(c_s-1\).

En el bosque la expansión conexa centrada de \pruebaen{P04} cancela las contracciones
gaussianas del único ciclo ordinario contra el medio logaritmo del determinante.
La única interacción excepcional que aún podría alcanzar \(w\) grupos contiene dos
inserciones con la misma amplitud. Podar una rama sustituye sus hojas por sus medias y
una derivada en el campo del grupo de enlace en ambas expresiones. En la raíz, con las
ocupaciones binarias \(Y_1,Y_2\), \pruebaen{X02} obtiene que el residuo tras la traza es
\(4(1-\mathbb E Y_1)(1-\mathbb E Y_2)\operatorname{Cov}(Y_1,Y_2)\).
Para alcanzar el soporte máximo esas fibras están en grupos distintos, independientes
también al variar sus campos. El residuo y sus derivadas se anulan. Quedan a lo sumo
\(w-1\) grupos, de modo que el primer coste es al menos \(c_s+1\).

Si hay un solo ciclo, para alcanzar el coste \(c_s-1\) todas las ramas fuera de su
componente se podan del mismo modo. Si \(k_i\) es el número de factores incidentes,
la razón ocupado/vacío en el producto de creencias es
\(x_i^{k_i}/[\prod_{f\ni i}q_f\,(x_i/(1-x_i))^{k_i-1}]=\zeta_i\) por estación.
El peso basal es \(Z_{\rm Bethe}^{-1}\). Las transiciones binarias del ciclo dan
\(Z/Z_{\rm Bethe}=1+\theta\). En el complemento del modo nulo axial, el lema
del determinante de \pruebaen{U07} reduce el cofactor a
\(\det[(-H_B''-R'')V_0]\prod_f\bigl(q_f/\prod_{i\in f}(1-x_i)\bigr)\), con el producto
de cocientes tomado sobre todas las fibras y \(V_0=\diag(x_i(1-x_i))\).
A reacción exterior cero, la identidad de incidencias de \pruebaen{Z01} da
\((1-\theta)^2\), un factor por orientación. El error es \(\log(1-\theta^2)\),
con \(\theta=(-1)^m\prod_{i\in\mathcal C}x_i/(1-x_i)\) para sus \(m\) sitios.
Cada \(x_i\) es divisible por \(\zeta_i\). Las derivadas de campo
\(\zeta_i\partial_{\zeta_i}\) conservan la doble multiplicidad. \pruebaen{Z02} y
\pruebaen{P05} dan así un campo adicional de cada sitio del ciclo en el término crítico.
Cuando hay dos o más ciclos basta la cuenta general de conexidad.

El paso al alcance creciente es \pruebaen{X03} y \pruebaen{V07}. Cambiar el orden alto
por \(q=4v+12\) conserva todos los grados leídos, ya que el cambio de germen está en
\((s^{q-1})\). Por identidad holomorfa se prolongan a toda la caja las anulaciones y
las dobles divisibilidades probadas inicialmente para cada soporte. Cauchy en \(s\)
desde el primer coste \(w\) da \(C^v(Cv/j)^w\). Schwarz sucesivo en activaciones
aporta \(\prod_i(Cd_i/j)\). En el coste crítico unicíclico los campos adicionales
cuestan a lo sumo \(\sum_i Cd_i/j\le Cv^2/j\) en el punto físico.
Así salen los tres exponentes con la constante anunciada. El tamaño del permanente
físico sólo interviene después, en la comparación real de la proposición~\ref{prop:m}.
\end{proof}
Para \(v\le4\) se usa la versión física a grado fijo de \pruebaen{P05}, incluida su
proyección del resto para cada tamaño entero. Los casos geométricos se especifican
en B.1 y en la pieza V05 del archivo de fuentes.

\section{Norma del error y cola de grados altos}
\begin{proposicion}\label{prop:n}
Sea \(P\) el punto estacionario de torres de la sección 1. Hay un umbral \(\eta>0\) de
clase tal que para \(\rho\ge2\) fijo, \(\rho^2(\epsilon_A+1/j)\le\eta\) y \(j\)
suficientemente grande, cada \(Z(Y)>0\) y
\[
 \|e\|_\rho:=\sum_{U\subseteq A}\rho^{|U|}|\Delta_Ue|\le C_\rho j^2.
\]
Para \(N=\lceil4\log n/\log\rho\rceil\) se tiene
\(\sum_{|U|>N}|\Delta_Ue|\le C_\rho n^{-2}\). Para \(n\) grande esta cota se
conserva al sustraer los errores de los cuadrados de B.1.
\end{proposicion}
\begin{proof}
En el álgebra finita \(z^U z^V=z^{U\cup V}\), la norma ponderada es submultiplicativa.
Sea \(a_U=\Delta_U{\bf1}_{\{\text{sin ataques diagonales}\}}\) y sea \(p_U\) la
probabilidad de prescribir \(U\) en la ley de torres, cero en caso de
incompatibilidad axial. La inversión booleana da \(Z=\sum_U a_U p_U z^U\)
antes de tomar logaritmos. En \pruebaen{V02}, para prescripciones cortas,
\(p_U=(\prod_{i\in U}b_i)e^{L_s(T_U)}(1+O(e^{-cj}))\),
con \(T_U\) la suma de las supresiones unitarias, \(\|T_U\|\le3|U|\).
La cota de Hessiana de \(L_s\) es \(C/j\) y los pesos \(b_i\) son comparables a \(D_i\).
Para todas las tallas se cumple \(p_U\le C_0^{|U|}\prod b_i\).

El polinomio de pesos \(w_U=a_Uz^U\prod b_i\) factoriza en componentes diagonales
y anula los sitios aislados. Un árbol etiquetado enraizado en \(i\) acota el total
de las componentes de tamaño \(v\) por
\(b_i(C_B\rho)^v v^{v-2}(C\epsilon_A)^{v-1}/(v-1)!\),
incluida una reserva \(e^{Bv}\) (\pruebaen{F04}).
Con \(S=j\rho^2(\epsilon_A+1/j)\) la masa total es \(O(S)\) y la de incompatibles
con una componente de tamaño \(v\) es \(O(vS/j)\). \pruebaen{V03b} demuestra por árboles
la convergencia absoluta del logaritmo del gas en el álgebra, incluso con la reserva.

En el gas inclinado por \(\exp(\xi[T_U])\), la media \(m_\xi\) de \(T_U\) y su derivada
tienen norma \(O(S)\) por las dos marcas de \pruebaen{V03}. La ecuación
\(\xi=DL_s(m_\xi)\) es contractiva con constante \(CS/j<1/2\).
Sea \(D_\xi=Z_\xi\exp(L_s(m_\xi)-\xi[m_\xi])\). Al centrar en \(m_\xi\) se
anula el término lineal. La generatriz de momentos centrados da la cota
\([C(\sqrt{\ell S}+\ell)]^\ell\) en grado \(\ell\). El Taylor del factor residual
tiene coeficientes de norma \((C/\sqrt{j\ell})^\ell\) hasta \(\ell=\lfloor\kappa j\rfloor\).
Su contracción desde grado dos cuesta \(C(1+S)/j\). La reserva exponencial controla
las prescripciones largas y el resto de Taylor, según las estimaciones de \pruebaen{V03}.
Así \(\|D_\xi^{-1} Z-1\|_\rho\le C(1+S)/j+Ce^{-c'j}<1/2\). En las evaluaciones
booleanas \(D_\xi>0\) por construcción real y \(Z(Y)>0\). La serie logarítmica da
\(\|F\|_\rho\le C(1+S)\). En particular \(Q(A)>0\) sobre \(H_j\).

Para \(\Lambda\), \pruebaen{N02} construye simultáneamente la silla de todas las caras
mediante el balanceo de \(P\) en el álgebra. En el calibre que elimina el modo nulo de
filas y columnas, la brecha da un inverso uniformemente acotado en norma de suma
por filas. La carga pequeña garantiza la contracción y \(q_e/P_e=1+O(\eta)\) en la
norma. La tangente en la rama pequeña identifica el máximo real en cada cara.
En \pruebaen{N03} la perturbación del complemento de Schur diagonal tiene norma
máxima de suma por filas \(<1/2\) y suma absoluta de entradas \(O(\rho^2j)\).
La serie de trazas cuesta \(O(\rho^2j)\), independientemente del número de diagonales.
La acción se acota usando \(\sum_e P_e|\log P_e|\le j\log j\), por entropía de filas.
Con los factores locales resulta \(\|\Lambda\|_\rho\le C(j\log(ej)+\rho^2j)\).

Se suman las dos normas y se extrae \(\rho^{-N}\le n^{-4}\), con \(j\le n\),
como en \pruebaen{Z04}. El grado de la sustracción por cuadrados es a lo sumo cuatro.
Una vez \(N\ge4\), todos los coeficientes de la cola coinciden con los de \(e\).
\end{proof}
La norma y sus dos demostraciones son deterministas. También lo es la positividad
que permite \(H_j\subset F_b\). En \(F_b\setminus H_j\) B.1 recurre al coste de
la línea larga y a la densidad detenida. Las estimaciones sobre la máscara corta
se aplican únicamente en \(H_j\).

\clearpage
\section{Demostraciones de balanceo, mezcla y átomos}
\pieza{D04}
\pieza{U02}
\pieza{V01}
\pieza{U03}
\pieza{U04}
\pieza{U07}
\pieza{A03}
\clearpage
\section{Demostraciones de las cancelaciones y de sus cotas uniformes}
\pieza{B01}
\pieza{B02}
\pieza{G02}
\pieza{G03}
\pieza{G04}
\pieza{P03}
\pieza{P04}
\pieza{X02}
\pieza{Z01}
\pieza{Z02}
\pieza{P05}
\pieza{X03}
\pieza{V07}
\clearpage
\section{Demostraciones de las dos normas de activaciones}
\pieza{V02}
\pieza{F04}
\pieza{V03}
\begingroup
\subsection{Convergencia de la expansión en polímeros}\label{pieza:V03b}
\noindent{\small\textsf{Fuente V03b}}\par\nobreak\smallskip
\paragraph{Justificación del criterio de polímeros usado en la parte 3.}
Sean $a(P)\ge0$ las normas de las actividades con el margen exponencial ya incluido.
La incompatibilidad incluye dos ocurrencias del mismo polímero. Por el argumento
anclado de F04 se cumple $\sum_{D\sim P}a(D)e^{|D|}\le |P|$ con holgura,
y $\sum_P a(P)e^{|P|}\le CS$, reservando también esta unidad en la exponencial.
Para comprobar la cota del logaritmo, introducir una actividad escalar auxiliar
y desarrollar primero formalmente. La función exponencial de los grafos conexos
enumera todos los grafos en las ocurrencias etiquetadas. Para una colección de $q$
ocurrencias, la suma sobre todos los grafos, con peso $-1$ por arista incompatible,
es el indicador de compatibilidad total por expansión del producto.
Por tanto el logaritmo retiene los grafos conexos y el factor $1/q!$.

En un grafo de incompatibilidad con sus aristas totalmente ordenadas, a cada subgrafo
conexo asígnese su árbol de Kruskal. Los subgrafos que dan un árbol $T$ se obtienen
añadiendo cualquier subconjunto de las aristas posteriores a todas las del camino
que une sus extremos en $T$: al procesarlas sus extremos ya están unidos.
Recíprocamente, cualquier arista extra anterior a la máxima del camino cambiaría
el árbol mínimo. La suma de signos sobre este intervalo es cero o tiene módulo uno.
La suma absoluta después de esta cancelación se domina por la suma sobre árboles.
Para árboles de altura a lo sumo $h$ con raíz $P$ y sin el peso de la raíz
el mayorante etiquetado exponencial se itera desde $F_0(P)=1$ mediante
\[
 F_{h+1}(P)=\exp\Big(\sum_{D\sim P}a(D)F_h(D)\Big).
\]
El factorial de los hijos proviene del orden irrelevante de las ramas; con raíz fijada
la serie lleva $1/(q-1)!$, frente a $1/q!$ en la serie sin raíz. Para $q=1$ el árbol es vacío de aristas.
Por inducción $F_h(P)\le e^{|P|}$. La suma sin raíz queda dominada por
$\sum_P a(P)e^{|P|}$ por designación de una raíz. El margen da convergencia
absoluta en un disco escalar mayor que la unidad: la desigualdad se conserva para
$(1+\delta')a(P)$ con $\delta'>0$ fijo al reducir $\eta_F$. Esto vale con actividades en
el álgebra booleana. La exponencial de la serie conexa coincide allí con la
partición por identidad de coeficientes. El margen restante en talla permite
marcar una o dos ocurrencias y efectuar el centrado de la parte 3.

\endgroup
\pieza{N02}
\pieza{N03}
\pieza{Z04}
\end{document}
